% Въ лѣто Господне ҂вк҃ѕ.
%
% Господи Іисусе Христе, Сыне Божій, помилуй меня грѣшнаго.

\documentclass[20pt]{extarticle}

\usepackage[a5paper, portrait, margin=6pt, top=0pt, bottom=0pt]{geometry}

%\usepackage{pgfpages}
%\pgfpagesuselayout{2 on 1}[a4paper, landscape]

\usepackage{fontspec}
\usepackage{xkeyval}
\usepackage{microtype}

\usepackage{polyglossia}
\setmainlanguage{churchslavonic}

\usepackage{churchslavonic}
\usepackage{lettrine}
\usepackage{graphicx}  % for \scalebox

\usepackage{hyperref}
\hypersetup{
	pdftitle={Икона Богородицы Толгская},
	pdfauthor={Православная Церковь}
}


\makeatletter

\def\@texttop{\vfill}
\def\@textbottom{\vfill}

\makeatother


\newcommand{\cuL}[1]{\scalebox{1.5}{\cuKinovarColor#1}}
\newcommand{\cul}[1]{\mbox{\cuKinovarColor#1}}

\newfontfamily\noto{Noto Sans Symbols 2}
\newfontfamily\churchslavonicfont[Script=Cyrillic,Ligatures=TeX,HyphenChar=_]{Acathist}
%\newfontfamily\cyrillicfontsf{DejaVu Sans}

\hyphenpenalty=10000
\setlength\emergencystretch{4em}
\def\hyph{{\_}{}{}\penalty-10000}

%\linespread{1.2}
\setlength{\parindent}{1cm}
\setlength{\parskip}{8pt}

\newcommand{\sub}[1]{\vspace{3pt}{\centering\cuKinovar{#1}\par}}


\begin{document}

\pagestyle{empty}
\centering

\sub{\large{\noto 🙚} И҆ко́на бж҃їей мт҃ре то́лгскаѧ {\noto 🙘}}

\sub{Тропа̀рь, глас д҃:}
\cuL Дне́сь свѣ́тлѡ сїѧ́етъ на то́лгѣ \\
ѻ҆́бразъ тво́й, пречⷭ҇таѧ дѣ́во бцⷣе, \\
\hspace{3pt}и҆ ꙗ҆́кѡ со́лнце незаходи́мое, \\
всегда́ вѣ́рнымъ подаде́сѧ, \\
є҆го́же ви́дѣвъ на воздꙋ́сѣ, \\
неви́димѡ а҆́гг҃лы, ꙗ҆́кѡ ники́мъ, держи́ма, \\
\hspace{4pt}преѡсвѧще́нный є҆пи́скопъ гра́да росто́ва трі́ѳѡнъ \\
\hspace{-3pt}тече́ ко ꙗ҆вле́нномꙋ свѣтѧ́щемꙋсѧ столпꙋ̀ ѻ҆́гненнꙋ, \\
и҆ по вода́мъ, ꙗ҆́кѡ по сꙋ́хꙋ, пре́йде, \\
и҆ молѧ́шесѧ тѝ вѣ́рнѡ ѡ҆ па́ствѣ и҆ ѡ҆ лю́дехъ. \\
\hspace{-2pt}\cul{И҆} мы́, къ тебѣ́ притека́юще, зове́мъ: \\
прест҃а́ѧ дв҃о бцⷣе, \\
вѣ́рнѡ тѧ̀ сла́вѧщихъ спаса́й, \\
\hspace{-3pt}странꙋ̀ на́шꙋ, а҆рхїере́євъ \\
и҆ всѧ̑ рѡссі́йскїѧ наро́ды ѿ всѣ́хъ бѣ́дъ и҆збавлѧ́й \\
\hspace{-3pt}по вели́цѣй твое́й млⷭ҇ти.


\newpage


\sub{Конда̀к, глас и҃:}
\cuL Чꙋ́дное и҆ бг҃опрїѧ́тное свѣтоно́снагѡ зра́ка твоегѡ̀ \\
\vspace{-4pt}\hfill ꙗ҆вле́нїе, пречⷭ҇таѧ дв҃о, \\
\hspace{3pt}съ вѣ́рою взира́ющымъ быва́етъ извѣ́стное \\
\vspace{-4pt}\hfill ча́ѧнїе сп҃се́нїѧ, \\
столпо́мъ бо ѻ҆́гненнымъ ꙗ҆влѧ́ющи предста́тельство, \\
неꙋсы́пное и҆ те́плое за ны̀ къ бг҃ꙋ возсыла́еши \\
\vspace{-5pt}\hfill моле́нїе, \\
и҆́мже странꙋ̀ рѡссі́йскꙋю ѿ всѧ́кихъ бѣ́дъ и҆зба́ви, \\
\vspace{-2pt}\hfill мо́лимсѧ, \\
да ра́достными ѹ҆сты̀ вопїе́мъ тѝ вси́: \\
ра́дꙋйсѧ, предводи́тельнице вѣ́чнагѡ бл҃же́нства.

\sub{Велича́ние:}
\cuL Величе́мъ тѧ̀, / прест҃а́ѧ дв҃о, / бг҃оизбра́ннаѧ ѻ҆трокови́це, / и҆ чти́мъ ѻ҆́бразъ тво́й ст҃ы́й, / и҆́мже то́чиши и҆сцѣлє́нїѧ // всѣ́мъ \\
съ вѣ́рою притека́ющымъ.


%\newpage
%\newgeometry{margin=1cm}
%\renewfontfamily\churchslavonicfont[Script=Cyrillic,Ligatures=TeX,HyphenChar=_]{Triodion}
%\churchslavonicfont\fontsize{17}{21}\selectfont


\end{document}
